\documentclass[../main.tex]{subfiles}
\begin{document}
\chapter{3D Solutions to the Schr\"odinger Equation}
\section{TISE in \texorpdfstring{$\R^3$}{3D}}
The TISE in $\R^3$ is given by:
\begin{equation}
  -\frac{\hbar^2}{2m} \nabla^2 \chi(\vec{x}) + U(\vec{x})\chi(\vec{x}) = E\chi(\vec{x}) \label{TISE3DDef}
\end{equation}
We would like to study spherically symmetric potentials $U(\vec{x}) = U(r)$ and so it is convenient to consider spherical polar coordinates.
\begin{remark}[Recap]
  Recall from IA Vector Calculus that spherical polar coordinates are defined by:
  \[
    \vec{x}(r, \theta, \varphi) = r\begin{pmatrix}
    \sin \theta \cos \varphi \\
    \sin \theta \sin \varphi \\
    \cos \theta \\
    \end{pmatrix} = r \vec{e}_r
  \]
  where $r \geq 0$, $\theta \in [0, \pi]$ and $\varphi \in [0, 2\pi)$:
  \begin{itemize}
    \item $r$ is called the \textit{radial distance} and is the distance the origin to $\vec{x}$.
    \item $\theta$ is called the \textit{polar angle} and is the angle between the $z$-axis and $\vec{x}$ measured down from the $z$ axis towards the $x$-$y$ plane
    \item $\varphi$ is called the \textit{azimuth angle} and is the angle measured anticlockwise from the $x$-axis round to the $y$-axis of the projection of $\vec{x}$ onto the $x$-$y$ plane.
  \end{itemize}
  Recall also that when performing a change of variables in a volume integral, we have to introduce a factor of the \textit{Jacobian} $|J| = r^2 \sin \theta$:
  \[
    \int_{\R^3} f(\vec{x}) \d{V} = \int_{0}^{2\pi} \int_{0}^{\pi} \int_{0}^{\infty} F(r, \theta, \varphi) r^2 \sin \theta \d{r} \d{\theta} \d{\varphi}
  \]
  where $F(r, \theta, \varphi) = f(\vec{x}(r, \theta, \varphi))$.

  Finally, the Laplacian in spherical polar coordinates is given by:
  \begin{equation}
    \nabla^2 f = \frac{1}{r^2} \pderiv{}{r}\left(r^2 \pderiv{f}{r}\right) + \frac{1}{r^2 \sin \theta} \pderiv{}{\theta}\left(\pderiv{f}{\theta} \sin \theta\right) + \frac{1}{r^2 \sin^2 \theta} \pderiv[2]{f}{\varphi} \label{lapSpher}
  \end{equation}
\end{remark}
In IA Dynamics and Relativity, we saw that when we had a spherically symmetric potential, the rate of change of angular momentum $\vec{L} = \vec{x} \times \vec{p}$ was always 0, that is:
\[
  \deriv{\vec{L}}{t} = \deriv{}{t}(\vec{x} \times \vec{p}) = \vec{0}
\]
Since $\vec{L}$ is constant and $\vec{L} \cdot \vec{x} = 0$, we concluded that the motion must be restricted to the plane $\vec{L} \cdot \vec{x} = 0$.
On this plane, we could then solve in the radial direction, reducing the problem to 1D motion, where we can use $\vec{L}$ to determine the angular component.
We would like to try and formulate a quantum mechanical version of this.
\section{Angular Momentum Operator}
\subsection{Definition and Properties}
In 3D, the momentum operator we defined in 1D in \cref{opExamples} becomes:
\[
  \hat{p} = -i\hbar \nabla
\]
and the position operator becomes:
\[
  \hat{x} = \vec{x}
\]
\begin{definition}[Angular momentum operator]
  We define the \textit{angular momentum operator} as:
  \[
    \hat{L} = \hat{x} \times \hat{p} = -i\hbar\vec{x} \times \nabla
  \]
\end{definition}
In Cartesian coordinates, the components of this are:
\begin{align*}
  \hat{L}_i &= [\hat{x} \times \hat{p}]_i \\
            &= \levi_{i j k} \hat{x}_j \hat{p}_k \\
            &= -i\hbar \levi_{i j k} x_j \pderiv{}{x_k}
\end{align*}
\subsubsection{Properties}
\begin{enumerate}
  \item Since $\hat{L}_i = \levi_{i j k} \hat{x}_j \hat{p}_k$ and the operators $\hat{x}_j = x_j$ and $\hat{p}_k = -i\hbar \pderiv{}{x_k}$ are both Hermitian, for any state $\psi \in \mathcal{H}$ we have:
    \begin{align*}
      (\hat{L}_i)^{\dag}\psi &= \levi_{i j k}(\hat{x}_j \hat{p}_k)^{\dag}\psi \\
                             &= \levi_{i j k}\hat{p}^{\dag}_k \hat{x}^{\dag}_j\psi \text{ by \cref{opProp} \textbf{(ii)}} \\
                             &= \levi_{i j k}\hat{p}_k \hat{x}_j\psi \text{ as $\hat{x}_j$ and $\hat{p}_k$ are Hermitian} \\
                             &= \levi_{i j k}\left(-i\hbar \pderiv{}{x_k}\right)(x_j\psi) \\
                             &= \levi_{i j k} x_j \left(-i\hbar \pderiv{}{x_k}\right)\psi \text{ as $\levi_{i j k} = 0$ if $j = k$} \\
                             &= \levi_{i j k} \hat{x}_j \hat{p}_k \psi \\
                             &= \hat{L}_i\psi
    \end{align*}
    Thus $(\hat{L}_i)^{\dag} = \hat{L}_i$, so $\hat{L}_i$ is Hermitian, as expected of an operator associated with a physical observable.
  \item The components of the angular momentum operator do not commute with each other, that is, $[\hat{L}_i, \hat{L}_j] \neq 0$ for $i \neq j$.
    \begin{align*}
      [\hat{L}_i, \hat{L}_j] &= \levi_{i p q} \levi_{j r s} [\hat{x}_p \hat{p}_q, \hat{x}_r \hat{p}_s] \\
                             &= \levi_{i p q} \levi_{j r s} (\hat{x}_p[\hat{p}_q, \hat{x}_r \hat{p}_s] + [\hat{x}_p, \hat{x}_r, \hat{p}_s]\hat{p}_q) \\
                             &= \levi_{i p q} \levi_{j r s} (\hat{x}_p([\hat{p}_q, \hat{x}_r]\hat{p}_s + \hat{x}_r[\hat{p}_q, \hat{p}_s]) + ([\hat{x}_p, \hat{x}_r]\hat{p}_s + \hat{x}_r[\hat{x}_p, \hat{p}_s])\hat{p}_q)
    \end{align*}
    We then see that since partials commute $[\hat{p}_q, \hat{p}_s] = 0$ and since multiplication is commutative $[\hat{x}_p, \hat{x}_r] = 0$.
    Moreover, for any $\psi \in \mathcal{H}$:
    \begin{align*}
      [\hat{p}_q, \hat{x}_r]\psi &= \hat{p}_q \hat{x}_r \psi - \hat{x}_r \hat{p}_q \psi \\
                                 &= -i\hbar\pderiv{}{x_q}(x_r \psi) + i\hbar x_r \pderiv{\psi}{x_q} \\
                                 &= -i\hbar\delta_{q r} \psi - i\hbar x_r \pderiv{\psi}{x_q} + i\hbar x_r \pderiv{\psi}{x_q} \\
                                 &= -i\hbar \delta_{q r} \psi
    \end{align*}
    Thus we have obtained the \textit{canonical commutation relation}:
    \begin{equation}
      [\hat{p}_q, \hat{x}_r] = -i\hbar\delta_{q r}\hat{I} \label{CCR}
    \end{equation}
    Substituting back in, we obtain:
    \begin{align*}
      [\hat{L}_i, \hat{L}_j] &= \levi_{i p q} \levi_{j r s} (\hat{x}_p (-i\hbar \delta_{q r}\hat{p}_s) + \hat{x}_r(i\hbar\delta_{p s}\hat{p}_q)) \\
                             &= -\levi_{i p q} \levi_{j q s} i\hbar \hat{x}_p \hat{p}_s + \levi_{i p q}\levi_{j r p} i\hbar \hat{x}_r \hat{p}_q \\
                             &= -(\delta_{i s}\delta_{p j} - \delta_{i j}\delta_{p s})i\hbar \hat{x}_p \hat{p}_s + (\delta_{q j}\delta_{i r} - \delta_{q r}\delta_{i j})i\hbar \hat{x}_r \hat{p}_q \\
                             &= i\hbar(-\hat{x}_j\hat{p}_i + \delta_{i j}\hat{x}_p\hat{p}_p + \hat{x}_i\hat{p}_j - \delta_{i j}\hat{x}_r\hat{p}_r) \\
                             &= i\hbar(\hat{x}_i\hat{p}_j - \hat{x}_j\hat{p}_i)
    \end{align*}
    Finally, we see that:
    \begin{align*}
      \levi_{i j k} \hat{L}_k &= \levi_{i j k}\levi_{p q k} \hat{x}_p \hat{p}_q \\
                              &= (\delta_{i p}\delta_{j q} - \delta_{i q}\delta_{j p})\hat{x}_p \hat{p}_q \\
                              &= \hat{x}_i \hat{p}_j - \hat{x}_j \hat{p}_i
    \end{align*}
    giving the result $[\hat{L}_i, \hat{L}_j] = i\hbar\levi_{i j k}\hat{L}_k$.

    Since $\hat{L}_i$ and $\hat{L}_j$ do not commute for $i \neq j$, we cannot measure different components of the angular momentum simultaneously.
    This means that we cannot have a well defined angular momentum in all three directions simultaneously.
\end{enumerate}
\begin{definition}[Total angular momentum operator]
  We define the \textit{total angular momentum operator} as:
  \[
    \hat{L}^2 = \hat{L}^2_1 + \hat{L}^2_2 + \hat{L}^2_3
  \]
\end{definition}
This does commute with each component of $\hat{L}$:
\begin{align*}
  [\hat{L}^2, \hat{L}_i] &= [\hat{L}_j \hat{L}_j, \hat{L}_i] \\
                         &= \hat{L}_j[\hat{L}_j, \hat{L}_i] + [\hat{L}_j, \hat{L}_i] \hat{L}_j \\
                         &= \hat{L}_j(i\hbar\levi_{j i k}\hat{L}_k) + i\hbar\levi_{j i k}\hat{L}_k \hat{L}_j \\
                         &= i\hbar\levi_{j i k}(\hat{L}_j\hat{L}_k + \hat{L}_k \hat{L}_j) \\
                         &= 0 \text{ since } \hat{L}_j\hat{L}_k + \hat{L}_k \hat{L}_j \text{ is symmetric in $j$ and $k$}
\end{align*}
Thus by \cref{simDiag}, $\hat{L}^2$ and $\hat{L}_i$ are simultaneously diagonalisable and so we can find a joint set of eigenstates for $\hat{L}^2$ and $\hat{L}_i$.
\subsection{Finding a Simultaneous Basis}
To find such a set of eigenstates, we first wish to write $\hat{L}_i$ and $\hat{L}^2$ in spherical polar coordinates.
For convenience, we look specifically at $\hat{L}_3$ as it has the nicest expression in spherical polar coordinate.
Using the multivariate chain rule, we obtain the expressions:
\begin{align}
  \hat{L}_3 &= -i\hbar \pderiv{}{\varphi} \nonumber \\
  \hat{L}^2 &= -\frac{\hbar^2}{\sin^2\theta}\left[\sin \theta \pderiv{}{\theta}\left(\sin \theta \pderiv{}{\theta}\right) + \pderiv[2]{}{\varphi}\right] \label{L2Spher}
\end{align}
We define the eigenvalue for $\hat{L}_3$ to be $\hbar m$ where $m$ is the \textit{azimuthal quantum number} and let $\lambda$ be the eigenvalue for $\hat{L}^2$.
Thus, the simultaneous eigenfunctions of $\hat{L}_3$ and $\hat{L}^2$ satisfy the system:
\begin{align}
  -i\hbar \pderiv{}{\varphi} Y(\theta, \varphi) &= \hbar m Y(\theta, \varphi) \label{spher1} \\
  -\frac{\hbar^2}{\sin^2\theta}\left[\sin \theta \pderiv{}{\theta}\left(\sin \theta \pderiv{}{\theta}\right) + \pderiv[2]{}{\varphi}\right] Y(\theta, \varphi) &= \lambda Y(\theta, \varphi) \label{spher2}
\end{align}
We looking for a separable solution of the form:
\[
  y(\theta, \varphi) = \Theta(\theta)\Phi(\varphi)
\]
Substituting into \cref{spher1}, we obtain:
\begin{align*}
  -i\hbar \Theta(\theta) \deriv{X}{\varphi} &= \hbar m \Phi(\varphi)\Theta(\theta) \\
  \implies \deriv{\Phi}{\varphi} &= i m\Phi(\varphi) \\
  \implies \Phi(\varphi) &= e^{im\varphi}
\end{align*}
Given that $Y$ must be single-valued in $\R^3$, $X$ must be $2\pi$ periodic, thus $e^{2i\pi m} = 1$, i.e. $m \in \Z$.

So:
\[
  \pderiv[2]{}{\varphi} Y(\theta, \varphi) = \Theta(\theta) \pderiv[2]{}{\varphi}(e^{im\varphi}) = -m^2 e^{im\varphi} \Theta(\theta)
\]
Substituting $Y(\theta, \varphi) = \Theta(\theta) \Phi(\varphi)$ into \cref{spher2} we obtain:
\begin{align*}
  -\frac{\hbar^2 e^{im\varphi}}{\sin^2 \theta} \left[\sin \theta \pderiv{}{\theta}\left(\sin \theta \pderiv{\Theta}{\theta}\right)\right] - \frac{\hbar^2 m^2e^{im\varphi}}{\sin^2 \theta} \Theta(\theta) &= \lambda e^{im\varphi}\Theta(\theta) \\
  \iff \frac{1}{\sin \theta} \pderiv{}{\theta}\left(\sin \theta \pderiv{\Theta}{\theta}\right) - \frac{m^2}{\sin^2 \theta}\Theta(\theta) &= - \frac{\lambda}{\hbar^2}\Theta(\theta)
\end{align*}
If we let $z = \cos \theta$, then this becomes the \textit{associated Legendre equation}.
\begin{align*}
  \Theta(\theta) &= P_{\ell, m}(\cos \theta) \\
                 &= (\sin \theta)^{|m|} \deriv[|m|]{}{(\cos \theta)} P_{\ell}(\cos \theta)
\end{align*}
where $P_{\ell, m}$ is the associated Legendre polynomial and $P_{\ell}$ is the ordinary Legendre polynomial.

We know that $P_{\ell}$ is a polynomial of degree $\ell$, and so for $Y$ to be non-zero, we must have $|m| \leq \ell$.
We call $\ell$ the \textit{total angular momentum quantum number}.
From the study of Legendre polynomials in IB Methods, we conclude that the eigenvalues of $\hat{L}^2$ are given by:
\[
  \lambda = \hbar^2 \ell(\ell + 1)
\]
So, putting everything together we obtain
\[
  Y_{\ell, m}(\theta, \varphi) = P_{\ell, m}(\cos \theta)e^{im\varphi}
\]
for $\ell \in \{0, 1, \ldots\}$ and $|m| \leq \ell$, satisfying:
\begin{align*}
  \hat{L}^2 Y_{\ell, m}(\theta, \varphi) &= \hbar^2 \ell(\ell + 1)Y_{\ell, m}(\theta, \varphi)  \\
  \hat{L}_3 Y_{\ell, m}(\theta, \varphi) &= \hbar m Y_{\ell, m}(\theta, \varphi)
\end{align*}
Thus $\{Y_{\ell, m}\}$ for $\ell \in \N_0$ and $|m| \leq \ell$ is a complete basis of $\mathcal{H}$ consisting of eigenfunctions for both $\hat{L}^2$ and $\hat{L}_3$.
\section{3D TISE for Spherically Symmetric Potentials}
\subsection{Relation to the Hamiltonian}
Comparing $\hat{L}^2$ in spherical polar coordinates (\cref{L2Spher}) and the Laplacian in spherical polar coordinates (\cref{lapSpher}), we observe that:
\begin{align*}
  \frac{\hat{L}^2}{r^2} &= -\frac{\hbar^2}{r^2 \sin^2 \theta}\left[\sin \theta \pderiv{}{\theta}\left(\sin \theta \pderiv{}{\theta}\right) + \pderiv[2]{}{\varphi}\right] \\
                        &= -\hbar^2 \left[\frac{1}{r^2} \pderiv{}{r}\left(r^2 \pderiv{}{r}\right) + \frac{1}{r^2 \sin^2 \theta}\left[\sin \theta \pderiv{}{\theta}\left(\sin \theta \pderiv{}{\theta}\right) + \pderiv[2]{}{\varphi}\right]\right]  + \frac{\hbar^2}{r^2} \pderiv{}{r}\left(r^2 \pderiv{}{r}\right) \\
                        &= -\hbar^2 \nabla^2 + \frac{\hbar^2}{r^2} \pderiv{}{r}\left(r^2 \pderiv{}{r}\right)
\end{align*}
From \cref{TISE3DDef}, the Hamiltonian $\hat{H}$ in $\R^3$ is:
\[
  \hat{H} = -\frac{\hbar^2}{2m} \nabla^2 + U(\vec{x})
\]
Thus, $\hat{H}$ in spherical polars for a spherically symmetric potential $U(\vec{x}) = U(r)$ can be written as:
\[
  \hat{H} = - \frac{\hbar^2}{2mr^2}\pderiv{}{r}\left(r^2 \pderiv{}{r}\right) + \frac{\hat{L}^2}{2mr^2} + U(r)
\]
Since $\hat{\vec{p}} = -i\hbar \nabla$, we can also write $\hat{\vec{p}}^2 = - \hbar^2 \nabla^2$ and so:
\begin{equation}
  \hat{H} = -\frac{1}{2m} \hat{\vec{p}}^2 + U(r) \label{hamilP}
\end{equation}
\subsubsection{Commutator Relations}
We would like to see if $\hat{H}$ commutes with $\hat{L}^2$ and $\hat{L}_3$.

For $\hat{H}$ and $\hat{L}_i$, we want to see if $\hat{L}_i$ commutes with $U(r) = U(|\vec{x}|)$.
Firstly:
\begin{align*}
  [\hat{L}_i, \hat{x}_j] &= [\levi_{i m n} \hat{x}_m \hat{p}_n, x_j] \\
                         &= \levi_{i m n} (\hat{x}_m [\hat{p}_n, \hat{x}_j] + [\hat{x}_m, \hat{x}_j]\hat{p}_n) \\
                         &= \levi_{i m n} \hat{x}_m [\hat{p}_n, \hat{x}_j] \text{ as $[\hat{x}_m, \hat{x}_j] = 0$} \\
                         &= \levi_{i m n} \hat{x}_m (-i\hbar \delta_{n j}) \text{ using \cref{CCR}} \\
                         &= i\hbar \levi_{i j m} \hat{x}_m
\end{align*}
We can then use this to show that:
\begin{align*}
  [\hat{L}_i, \hat{x}^2_j] &= [\hat{L}_i, \hat{x}_j]\hat{x}_j + \hat{x}_j[\hat{L}_i, \hat{x}_j] \\
                           &= i\hbar \levi_{i j m} (\hat{x}_m \hat{x}_j + \hat{x}_j \hat{x}_m) \\
                           &= 0 \text{ as $\hat{x}_m \hat{x}_j + \hat{x}_j \hat{x}_m$ is symmetric in $j$ and $m$}
\end{align*}
The operator which measures $|\vec{x}|^2$ is $\hat{\vec{x}}^2$, which therefore has:
\begin{align*}
  [\hat{L}_i, \hat{\vec{x}}^2] &= [\hat{L}_i, \hat{x}^2_1 + \hat{x}^2_2, + \hat{x}^2_3]  \\
                               &= [\hat{L}_i, \hat{x}^2_1] + [\hat{L}_i, \hat{x}^2_2] + [\hat{L}_i, \hat{x}^2_3] \\
                               &= 0
\end{align*}
Thus since $r = |\vec{x}|$ and $U(r)$ are unaffected by rotation:
\[
  [\hat{L}_i, U(r)] = \left[\hat{L}_i, U\left(\sqrt{\hat{\vec{x}}^2}\right)\right] = 0
\]

We also need to know if $\hat{L}_i$ commutes with $\vec{p}^2$.
Firstly:
\begin{align*}
  [\hat{L}_i, \hat{p}_j] &= [\levi_{i m n} \hat{x}_m \hat{p}_n, \hat{p}_j] \\
                         &= \levi_{i m n} (\hat{x}_m [\hat{p}_n, \hat{p}_j] + [\hat{x}_m, \hat{p}_j]\hat{p}_n) \\
                         &= \levi_{i m n} \hat{p}_n [\hat{x}_m, \hat{p}_j] \text{ as $[\hat{p}_n, \hat{p}_j] = 0$} \\
                         &= \levi_{i m n} \hat{p}_n (i\hbar \delta_{m j}) \text{ using \cref{CCR}} \\
                         &= i\hbar \levi_{i j n} \hat{p}_n
\end{align*}
Similarly to how we dealt with $\hat{x}^2_j$, we have:
\begin{align*}
  [\hat{L}_i, \hat{p}^2_j] &= [\hat{L}_i, \hat{p}_j]\hat{p}_j + \hat{p}_j[\hat{L}_i, \hat{p}_j] \\
                           &= i\hbar \levi_{i j n} (\hat{p}_n\hat{p}_j + \hat{p}_j \hat{p}_n) \\
                           &= 0 \text{ as $\hat{p}_n \hat{p}_j + \hat{p}_j \hat{p}_n$ is symmetric in $j$ and $n$}
\end{align*}
So by the same argument as for $\hat{\vec{x}}^2$, we see that:
\[
  [\hat{L}_i, \hat{\vec{p}}^2] = 0
\]
Combing both the results for $\hat{\vec{p}}$ and $\hat{\vec{x}}$, using the expression for $\hat{H}$ in \cref{hamilP}, we have:
\begin{align*}
  [\hat{L}_i, \hat{H}] &= \left[\hat{L}_i, -\frac{1}{2m} \hat{\vec{p}}^2 + U(r)\right] \\
                       &= -\frac{1}{2m} [\hat{L}_i, \hat{\vec{p}}^2] +  [\hat{L}_i, U(r)] \\
                       &= 0
\end{align*}
So $\hat{L}_i$ and $\hat{H}$ commute.
From this, we know that $[\hat{L}^2_i, \hat{H}] = 0$, and so using $\hat{L}^2 = \hat{L}^2_1 + \hat{L}^2_2 + \hat{L}^2_3$, we have also that $[\hat{L}^2, \hat{H}] = 0$.

We therefore have that $\{\hat{H}, \hat{L}^2, \hat{L}_3\}$ is a set of mutually commuting operators.
This means that we can find a complete basis for $\mathcal{H}$ which consists of states which are eigenstates for all three of these operators.
In terms of measurements, this means that we can measure $|\hat{L}|$, $L_z$ and $E$ simultaneously up to arbitrary precision.
\begin{remark}
  It can be shown that this set of operators is \textit{maximal}, that is, we cannot construct another independent Hermitian operator (other than $\hat{I}$) that commutes with all three operators.
\end{remark}
\subsection{Finding Spherically Symmetric Eigenfunctions}
Taking $\ell = 0$ forces $m = 0$ which gives $Y_{0, 0} = \Theta(\theta) \Phi(\varphi) = C$ for some normalisation constant $C$ and eigenvalue $\lambda = 0$.
So in this case, $\chi(r, \theta, \varphi) = Y_{0, 0}(\theta, \varphi) R(r) = C R(r)$, i.e. we are looking for spherically symmetric eigenfunctions.
Since $\hat{L}_3$ has no $r$ derivatives or constant terms, $\chi$ will always be an eigenstate of $\hat{L}_3$.
For $\chi = C R(r)$ to be an eigenstate of $\hat{L}^2$, we require:
\[
  \hat{L}^2 R(r) = ER(r)
\]
which gives:
\begin{align*}
  -\frac{h^2}{2mr}\deriv[2]{}{r}\left(r R(r)\right) + U(r) R(r) &= ER(r) \\
  \iff -\frac{h^2}{2m}\left(\deriv[2]{R}{r} + \frac{2}{r} \deriv{R}{r}\right) + U(r) R(r) &= ER(r)
\end{align*}
For $\chi$ to be normalisable, we need:
\begin{align*}
  \int_{\R^3} |\chi|^2 \d{V} = \int_{0}^{2\pi} \int_{0}^{\pi} \int_{0}^{\infty} R(r)^2 r^2 \sin \theta \d{r} \d{\theta} &= 1 \\
  \implies 4\pi \int_{0}^{\infty} r^2R(r)^2 \d{r} &= 1
\end{align*}
For the above integral to converge, we need that $R(r) \to 0$ sufficiently fast as $r \to \infty$.

Instead of using the above, it is instead more convenient to define $\sigma(r) = r R(r)$.
We then have that $\sigma(r)$ needs to satisfy:
\[
  -\frac{\hbar^2}{2m} \pderiv[2]{\sigma}{r} + U(r)\sigma(r) = E \sigma(r)
\]
This is exactly the 1D TISE (\cref{1DTISE}) but along the radial direction.
Since $r \in \R^{+}$, the normalisation condition for $\sigma$ is:
\[
  \int_{0}^{\infty} \sigma(r)^2 \d{r} = \frac{1}{4 \pi}
\]
We decide that we want to find a solution with $\sigma(0) = 0$ as imposing the regularity condition $\sigma(0) = a \neq 0$ leads to a $\chi$ which does not belong to the domain of a physical (i.e. Hermitian) Hamiltonian.

Continued next lecture.
\end{document}
